% =====================================================================
%  Supplementary Note (DRAFT, not yet \input by supplement.tex):
%  credibility tiers, why UMLS has no proxy, and the external alignment.
%  Generated from the credibility sweep (protocol data/umls/sweeps/
%  credibility.yaml; results are not redistributed -- regenerate with
%  gem-umls-sweep + gem-umls-classify-credibility against a licensed UMLS
%  copy) and data/umls/adjudications.yaml (credibility/*). Review before use.
% =====================================================================
\section{Credibility tiers: definitions, the UMLS finding, and external alignment}
\label{sec:supp-credibility}

\subsection{What the tiers mean}
\label{sec:supp-credibility-defs}

The overall credibility rating is the only subjective core dimension of the
model. Its four levels are defined by \emph{defeasibility} --- what it would
take for the curator to stop believing the evidence item --- rather than by
study design or by the authors' expressed certainty:

\begin{description}
  \item[\dimval{VERY\_HIGH}] believed by default even when directly
    contradicted; only an opposing argument of greater strength prompts
    re-examination;
  \item[\dimval{HIGH}] believed unless directly contradicted by evidence of
    comparable strength; such a contradiction opens doubt;
  \item[\dimval{MEDIUM}] not believed on its own; believed once corroborated
    by an independent source or an orthogonal evidence type;
  \item[\dimval{LOW}] not relied upon; recorded pending independent
    replication.
\end{description}

The definition is a belief-revision policy, which is what SEPIO's confidence
attribute --- the degree of belief that the asserted proposition is
true~\cite{sepio} --- asks for; the facets (cohort size, replication,
multiple-testing control, ancestry control, ascertainment) are the inputs to
that judgment and are kept separate, following the GRADE decomposition
principle~\cite{grade2008}. The scale intensifies at the top, as ACMG
evidence strength (\emph{supporting} to \emph{very strong})~\cite{acmg2015}
and ClinGen gene--disease validity (\emph{limited} to
\emph{definitive})~\cite{clingen2018} do; GRADE certainty intensifies at
the bottom (\emph{very low}) because it grades a body of evidence by
downgrading from a design-determined starting point. The two constructs are
therefore siblings, and the alignment in \cref{tab:cred-alignment} is
deliberately non-positional.

\subsection{Why the tiers have no UMLS proxy}
\label{sec:supp-credibility-umls}

An \emph{unmapped} entry in the UMLS crosswalk is a claim, not an absence,
so the credibility tiers were tested by a pre-registered sweep
(\texttt{data/umls/sweeps/credibility.yaml}). Adequacy criteria were fixed
before searching: a candidate concept must (A) denote a degree of epistemic
warrant rather than the magnitude of a measured property; (B) be a scale
point rather than the dimension; (C) be usable in a set that is disjoint and
ordered as the GEM scale is; and (D) carry that meaning in the relevant
domain rather than as a homonym. Three passes ran against a local index of
UMLS 2026AA (%
174 queries; 497 distinct concepts; replicated
on the UTS REST API, 505 concepts): whole-string matches for
intensified and degree qualifiers, the names of evidence-grading
vocabularies, and the same terms searched within the semantic types that
host scale values (Qualitative, Quantitative and Idea-or-Concept, Intellectual
Product, Finding). Every concept surfaced was assigned the criterion it
fails (\cref{tab:cred-families}; the full table is generated from the
results file).

\begin{table}[H]
  \centering
  \small
  \caption[Credibility sweep: rejected families]{Concepts surfaced by the
  credibility sweep (UMLS 2026AA), grouped by the reason they fail the
  adequacy criteria. Criterion codes as in the text.}
  \label{tab:cred-families}
  \begin{tabularx}{\linewidth}{@{}>{\raggedright\arraybackslash}X r c >{\raggedright\arraybackslash}X@{}}
    \toprule
    Family & $n$ & Fails & Examples \\
    \midrule
    'grade' (tumour/school/product) & 202 & D & Astrocytoma, low grade (C1314694) \\
    LOINC/PROMIS self-efficacy items & 191 & D & Current level of confidence I can drive a car (C5142440) \\
    SNOMED qualifier values (degree) & 21 & A & Very high (C0442804); Extremely high (C5787630) \\
    clinical 'evidence of' (sign present) & 20 & D & Blood alcohol level of less than 20 mg/100 ml (C0481372); Blood alcohol level of 20-39 mg/100 ml (C0481373) \\
    qualifier used as a measurand degree & 12 & A & Observation Value - High (C1561957); Message Waiting Priority - High (C1561958) \\
    scores, probabilities, risk categories & 8 & A & IPSS-R Risk Category Very High (C5202916); IPSS Risk Category High (C4522209) \\
    patient-perceived credibility of information & 7 & D & Establish teacher credibility, as appropriate (C0511083); credibility (C0870373) \\
    NCI/PDQ Level of Evidence I–IV & 6 & A & Level of Evidence (C0393009); Level of Evidence I (C2986612) \\
    metabolite identification confidence & 6 & D & Level 1 Metabolite Identification Confidence (C5551417); Level 3 Metabolite Identification Confidence (C5551419) \\
    SNOMED diagnostic-certainty modality & 5 & C & Definite (C0439544); Probable diagnosis (C0332148) \\
    NCI Confidence answer set & 4 & D & Very High Confidence (C4331500); High Confidence (C4330261) \\
    self-confidence (psychological) & 2 & D & Euphoric mood (C0235146); Confidence level (self-esteem) (C0518578) \\
    lexical collisions (media, lower) & 2 & D & Communications Media (C0009458); Culture Media (C0010454) \\
    degree of a measured property & 2 & A & A Medium Amount of Time (C4522282); A Medium Amount (C4522283) \\
    AE causality ladder (NCI) & 2 & D & Definitely Related to Intervention (C1704787); Possibly Related to Intervention (C1705910) \\
    LOINC molecular-pathology level of evidence & 2 & D & Level of evidence:Find:Pt:Bld/Tiss:Nom (C4760312) \\
    diagnostic-certainty rating items & 2 & D & UHDRS 1999 Version - Diagnosis Confidence Level (C4744064); Level of Diagnostic Certainty (C5909371) \\
    TNM certainty factor & 2 & D & Degree of certainty of TNM classification (C0456878); Qualifier for TNM degree of certainty (C1302640) \\
    statistical confidence quantities & 1 & A & Level of Confidence (statistical) (C5702556) \\
    \bottomrule
  \end{tabularx}
\end{table}

Four families deserve explicit argument because they are the closest
constructs UMLS holds. (i) The SNOMED qualifier values
\emph{High}~(C0205250, SNOMED 75540009), \emph{Very high}~(C0442804),
\emph{Low}~(C0205251) and their kin sit under
\emph{Increased}/\emph{Decreased}: the source terminology itself classifies
them as magnitudes of a measurand, a different relation from a degree of
warrant (A). (ii) The
NCI Thesaurus \emph{Level of Evidence I--IV} is a study-design hierarchy;
mapping credibility to it would fold the model's \emph{method} dimension into
credibility, contrary to the decomposition the model adopts (A, and C for the
scale). (iii) The NCI \emph{Very High / High / Moderate / Low Confidence}
set has exactly the shape of the GEM scale --- four levels, intensified at
the top --- but is a \emph{Clinical or Research Assessment Answer}: a
questionnaire response option about a respondent's confidence in performing
an activity (D). (iv) The SNOMED diagnostic-certainty modality
\emph{definite / probable / possible} is the nearest epistemic construct: it
denotes the likelihood that a proposition holds, but it is a three-point
ladder with no intensified top, bound in its sources to diagnosis
assertions, and it grades likelihood rather than defeasibility (C, D). The
same holds for the TNM certainty factor, the LOINC molecular-pathology
\emph{level of evidence} (an AMP/ASCO/CAP actionability tier for a variant
report) and the metabolite-identification confidence levels: ordered
epistemic ladders, each bound to a different object. Nor does a lexical
route --- post-coordinating a qualifier concept with an intensifier from the
SPECIALIST Lexicon (\texttt{LRMOD}) --- supply the missing degrees: the
lexicon's \emph{intensifier} class is a syntactic bucket (\emph{very},
\emph{slightly}, \emph{about}, \emph{ago}) that carries neither magnitude nor
direction, so a lexical identifier adds nothing a bare string does not; this
is an argument about what identifiers are for, not an empirical finding.

\subsection{External alignment}
\label{sec:supp-credibility-align}

\begin{table}[H]
  \centering
  \small
  \caption[Credibility: external alignment]{Non-positional alignment of the
  GEM credibility tiers with external scales. ``---'' marks a level with no
  counterpart; the FHIR column is the lossy export used when a
  \texttt{certainty-rating} code is required~\cite{fhir_evidence}.}
  \label{tab:cred-alignment}
  \begin{tabularx}{\linewidth}{@{}l >{\raggedright\arraybackslash}X >{\raggedright\arraybackslash}X >{\raggedright\arraybackslash}X >{\raggedright\arraybackslash}X >{\raggedright\arraybackslash}X@{}}
    \toprule
    GEM & GRADE certainty~\cite{grade2008} & FHIR \texttt{certainty-rating} & CIO confidence level & ACMG evidence strength~\cite{acmg2015} & ClinGen validity~\cite{clingen2018} \\
    \midrule
    \dimval{VERY\_HIGH} & --- (GRADE caps at \emph{high}) & \texttt{high} (lossy) & high (partial) & very strong & definitive \\
    \dimval{HIGH}       & high     & \texttt{high}     & high   & strong     & strong \\
    \dimval{MEDIUM}     & moderate & \texttt{moderate} & medium & moderate   & moderate \\
    \dimval{LOW}        & low      & \texttt{low}      & low (``absence of trust'', nearer to none) & supporting & limited \\
    ---                 & very low & \texttt{very-low} & ---    & ---        & --- \\
    \bottomrule
  \end{tabularx}
\end{table}

The alignments belong to the model, not to the UMLS crosswalk, which
records UMLS mappings or their considered absence.
